\mode*
\section{Introduction}

LaTeX is the standard tool for technical writing in large parts of science
and engineering, and many curricula introduce it early.
Unlike introductory programming, however, learners' difficulties with
LaTeX are barely researched.
The one controlled comparison we found is sobering: LaTeX users---experts
included---were slower and produced more typesetting, orthographical,
grammatical and formatting errors than Word users on continuous and
tabular text, although they reported enjoying their tool more, and
equation-heavy text was LaTeX's strong case \autocite{Knauff2014Latex}.
For teaching, this cuts two ways: the learning curve is real and must be
budgeted for, and the honest motivation for LaTeX is not general
efficiency but equations, structure, and versionable plain text.

\mode<presentation>{%
\begin{frame}
  \begin{block}{LaTeX and learners}
    \begin{itemize}
      \item Widely required, barely researched.
      \item Controlled comparison: slower, more errors than Word on plain
        text; stronger on equations; more enjoyment
        \autocite{Knauff2014Latex}.
      \item Honest motivation: equations, structure, plain text---not
        general speed.
    \end{itemize}
  \end{block}
\end{frame}%
}

The learner also arrives with a strong prior model: the WYSIWYG word
processor, where the document \emph{is} what the screen shows.
LaTeX breaks that model---source and output are distinct artifacts related
by compilation---much as the notional machine of programming breaks the
everyday model of the computer
\autocite{Sorva2013NotionalMachines,DuBoulay1981BlackBox}; difficulties
tied to inaccurate mental models are well documented for programming
\autocite{QianLehman2017Misconceptions,BenAri1998Constructivism}, and we
conjecture the same mechanism here.

In this work we apply the approach of our companion paper on introductory
programming \autocite{VTProgMisconceptions}: we treat documented
difficulties as the phenomenographic observations for a variation-theory
analysis \autocite{NCOL}, identify the critical aspects learners must
discern, and outline patterns of variation to teach them.
We cover the compile model, error handling, and semantic markup.

We ask the following research questions.

\begin{frame}
\begin{restatable}{question}{rqmisconceptions}\label{rq:misconceptions}
  Which misconceptions and difficulties with LaTeX---the compile model,
  error handling, and semantic markup---do the literature and our course
  data document?
\end{restatable}

\begin{restatable}{question}{rqaspects}\label{rq:aspects}
  Which critical aspects of document preparation with LaTeX do these
  misconceptions and difficulties point to?
\end{restatable}

\begin{restatable}{question}{rqpatterns}\label{rq:patterns}
  Which patterns of variation can make these critical aspects discernible
  to learners?
\end{restatable}
\end{frame}

Our contributions are:
a review of the (scarce) literature on learners' difficulties with LaTeX,
answering \cref{rq:misconceptions};
a variation-theory analysis identifying the critical aspects, answering
\cref{rq:aspects};
and a set of patterns of variation for teaching LaTeX, answering
\cref{rq:patterns}.
% TODO: Given the literature gap, this paper will need our own course data
% (datintro LaTeX lab grading and tutoring observations) more than the
% companions do; plan that data collection explicitly in the method.
